% Paper 5, §3 — Knowledge benchmarks: the MMLU lineage
% Anchors: kb/notes/evaluation/knowledge-benchmarks.md
%          kb/excerpts/{hendrycks2021-mmlu,mmlu-redux-2024,wang2024-mmlu-pro,
%                       rein2023-gpqa,glazer2024-frontiermath,phan2025-hle}.md

\section{Knowledge benchmarks: the MMLU lineage}
\label{sec:knowledge-benchmarks}

What we measure when we say a model \emph{knows things} is the answer
to a multi-subject multiple-choice exam, scored against a fixed answer
\glsterm{key}{key}, on the assumption that the answer \glsterm{key}{key} is right. The lineage that
made this concrete --- \glsterm{mmlu}{MMLU} \citep{hendrycks2021-mmlu}, its 2024
verifier audit \citep{mmlu-redux-2024}, the \glsterm{mmlu-pro}{MMLU-Pro} extension
\citep{wang2024-mmlu-pro}, \glsterm{gpqa-diamond-2}{GPQA Diamond} \citep{rein2023-gpqa},
\glsterm{frontiermath}{FrontierMath} \citep{glazer2024-frontiermath}, and \glsterm{hle}{Humanity's Last Exam}
\citep{phan2025-hle} --- is the central instrument of frontier
knowledge measurement from 2021 to 2026. It is also the lineage where
the field's reported numbers and the benchmark's actual discriminative
ceiling have most visibly come apart. The thesis of this section is
narrow: at the frontier, \glsterm{mmlu}{MMLU} is no longer a measurement instrument
but a noise-floor smoke test, and every successor has been built to
restore signal --- by raising difficulty, by widening the choice set,
by demanding reasoning rather than retrieval, or by deliberately
constructing items the field has not yet seen.

\subsection{MMLU: the canonical knowledge benchmark}
\label{sec:knowledge-mmlu}

\citet{hendrycks2021-mmlu} introduces \glsterm{mmlu}{MMLU} as a multi-subject test of
factual breadth and depth. The benchmark has $15{,}908$ four-way
multiple-choice questions across $57$ subjects --- elementary
mathematics, US history, computer science, law, medicine, ethics, and
others --- grouped into four broad areas (humanities, social sciences,
STEM, and other-professional)%
\footnote{KB excerpt: \texttt{kb/excerpts/hendrycks2021-mmlu.md\#sec-construction}.}.
Items are sourced from publicly available exam banks (AP exams, GRE,
MCAT, professional licensing exams) and filtered for clarity and
unique correct answer; the verifier $V(\hat{a}, a^\star) \in \{0, 1\}$
is exact-match on the chosen letter. Random-guess baseline is $25\%$;
the human-expert baseline cited in the paper is $\sim 89.8\%$
averaged across subjects \citep[\S{}abstract]{hendrycks2021-mmlu}.

The headline launch result is the load-bearing one
\citep[\S{}abstract]{hendrycks2021-mmlu}:

\begin{quote}\itshape
We find that while most recent models have near random-chance accuracy,
the very largest GPT-3 model improves over random chance by almost 20
percentage points on average. However, on every one of the 57 tasks,
the best models still need substantial improvements before they can
reach expert-level accuracy.
\end{quote}

GPT-3 175B at launch scored $\sim 43.9\%$, near random on many subjects,
strong on humanities and weaker on calculations and specialized
professional knowledge. This is the gap \glsterm{mmlu}{MMLU} was constructed to expose:
broad-but-shallow retrieval, not multi-step reasoning. Two operational
choices made \glsterm{mmlu}{MMLU} the headline benchmark of 2021--2024. First, the
EleutherAI \texttt{\glsterm{lm-evaluation-harness}{lm-evaluation-harness}} loglikelihood-scoring recipe
\citep[\S{}\ref{sec:knowledge-mmlu}]{hendrycks2021-mmlu} --- five-shot
in-context, direct-answer (no \glsterm{chain-of-thought}{chain-of-thought}), pick the argmax over
$\{\log P(\text{A}), \log P(\text{B}), \log P(\text{C}),
\log P(\text{D})\}$ --- made open-weight scores reproducible across
groups. Second, the $57$-subject decomposition gave a per-subject
diagnostic surface: a model could be evaluated not only on aggregate
accuracy but on subject-area lopsidedness, which is the construct the
abstract names explicitly. The combination is what made \glsterm{mmlu}{MMLU} \emph{the}
number on every model card from GPT-3 forward.

\subsection{The 6.49\% verifier-error ceiling}
\label{sec:knowledge-redux}

The first thing the field broke on \glsterm{mmlu}{MMLU} was the assumption that the
answer \glsterm{key}{key} was right. \citet{mmlu-redux-2024} (\emph{Are We Done with
\glsterm{mmlu}{MMLU}?}) systematically audit $5{,}700$ \glsterm{mmlu}{MMLU} questions against expert
annotators and report a load-bearing number%
\footnote{KB excerpt: \texttt{kb/excerpts/mmlu-redux-2024.md\#abstract}.}
\citep[\S{}abstract]{mmlu-redux-2024}:

\begin{quote}\itshape
6.49\% of \glsterm{mmlu}{MMLU} questions contain errors that range from incorrect
ground truth labels to ambiguous question text and answers.
\end{quote}

The errors include malformed questions, wrong gold answers, ambiguous
wording, and items where multiple options are defensibly correct. The
distribution is heavy-tailed: in the Virology subject, $57\%$ of the
analyzed questions contain errors, which means the per-subject score
on Virology is almost entirely verifier noise
\citep[\S{}abstract]{mmlu-redux-2024}.

The implication for headline scores is mechanical. If $6.49\%$ of the
gold labels are wrong, a perfect oracle scores at most
$\sim 93.5\%$ on gold-labeled \glsterm{mmlu}{MMLU}. Frontier models clustered at
$90\%\text{--}92\%$ in 2024--25 are \emph{inside} the verifier-error
band of one another and of the perfect-oracle ceiling --- leaderboard
rankings inside that band are unstable to verifier corrections. This
is not a complaint about the audit; it is what the audit found.

\begin{contradiction}
The contradiction is not between two papers but between what the field
\emph{reports} and what the benchmark can \emph{measure}. Vendors and
research papers regularly publish \glsterm{mmlu}{MMLU} scores at single-percentage-point
precision in the saturated regime: \(91.2\%\), \(92.4\%\), \(93.0\%\).
Per \citet{mmlu-redux-2024}, none of these numbers is meaningfully
distinguishable above the verifier-error band; any score reported
above $\sim 93.5\%$ is by construction unverifiable against a corrected
gold standard, and differences between two models reported at $90.5\%$
and $92.0\%$ are within audit noise. The field continues to publish
these numbers as if \glsterm{mmlu-redux-2}{MMLU-Redux} had not been written.
\end{contradiction}

The audit is a methodology contribution, not a new benchmark. Polo et
al.\ publish per-question corrections so any subsequent run can use
audited gold labels, but adoption has been partial: the Open \glsterm{llm}{LLM}
Leaderboard moved its primary headline to \glsterm{mmlu-pro}{MMLU-Pro}
(\cref{sec:knowledge-mmlu-pro}), while many model cards still report
raw \glsterm{mmlu}{MMLU} as an in-distribution sanity check.

\subsection{MMLU-Pro: the 24-prompt-style robustness extension}
\label{sec:knowledge-mmlu-pro}

\citet{wang2024-mmlu-pro} address \glsterm{mmlu}{MMLU}'s \glsterm{saturation}{saturation} directly with three
structural changes%
\footnote{KB excerpt: \texttt{kb/excerpts/wang2024-mmlu-pro.md\#abstract}.}
\citep[\S{}abstract]{wang2024-mmlu-pro}:

\begin{quote}\itshape
This paper introduces \glsterm{mmlu-pro}{MMLU-Pro}, an enhanced dataset designed to extend
the mostly knowledge-driven \glsterm{mmlu}{MMLU} benchmark by integrating more
challenging, reasoning-focused questions and expanding the choice set
from four to ten options. Additionally, \glsterm{mmlu-pro}{MMLU-Pro} eliminates the trivial
and noisy questions in \glsterm{mmlu}{MMLU}.
\end{quote}

The structural changes are: (i) $4$-way $\to$ $10$-way MCQ, dropping
the \glsterm{random-guess-ceiling}{random-guess ceiling} from $25\%$ to $10\%$ and expanding dynamic
range above the noise floor; (ii) reasoning-focused item authoring;
(iii) implicit denoising of the trivial / noisy items the
\citeauthor{mmlu-redux-2024} audit identified. At launch, top frontier
models scored $16$--$33$ percentage points lower on \glsterm{mmlu-pro}{MMLU-Pro} than on
\glsterm{mmlu}{MMLU} \citep[\S{}sec-delta]{wang2024-mmlu-pro}.

The most original contribution is the prompt-stability protocol. The
authors evaluate each model under $24$ different prompt styles and
report the resulting score variance%
\footnote{KB excerpt: \texttt{kb/excerpts/wang2024-mmlu-pro.md\#sec-delta}.}
\citep[\S{}abstract]{wang2024-mmlu-pro}:

\begin{quote}\itshape
With 24 different prompt styles tested, the sensitivity of model
scores to prompt variations decreased from 4--5\% in \glsterm{mmlu}{MMLU} to just
2\% in \glsterm{mmlu-pro}{MMLU-Pro}.
\end{quote}

A $4$--$5\%$ prompt-induced score variance on \glsterm{mmlu}{MMLU} is significant ---
it can reorder leaderboard positions on its own --- and the $10$-way
design halves it both by diluting the \glsterm{answer-letter-prior-2}{answer-letter prior} (a
pre-trained model's bias toward A/B/C/D contributes a smaller fraction
of the total score with more distractors) and by selecting items where
the answer is determined by the question content rather than the
surface format. The implication is that prompt-stability is now a
benchmark axis, not an evaluation hygiene afterthought: the same
model on the same items under different prompt templates can produce
scores that differ by more than the gap between two models.

A second diagnostic falls out of the \glsterm{mmlu-pro}{MMLU-Pro} construction. On the
original \glsterm{mmlu}{MMLU}, \glsterm{chain-of-thought}{chain-of-thought} and direct-answer scoring give
comparable accuracy; on \glsterm{mmlu-pro}{MMLU-Pro}, \glsterm{chain-of-thought}{CoT} gives a marked improvement
\citep[\S{}sec-cot]{wang2024-mmlu-pro}:

\begin{quote}\itshape
\glsterm{chain-of-thought}{Chain-of-Thought} reasoning showed marked improvements on \glsterm{mmlu-pro}{MMLU-Pro},
contrasting with the original \glsterm{mmlu}{MMLU} where direct answering performed
comparably---indicating the benchmark includes substantially more
complex reasoning questions.
\end{quote}

\begin{intuition}
A benchmark where \glsterm{chain-of-thought}{CoT}-vs-direct scoring gives the same number is
dominated by retrieval: spending inference compute on intermediate
reasoning does not help, because the items are answerable in one
step given the relevant fact. A benchmark where \glsterm{chain-of-thought}{CoT} measurably
outperforms direct is exposing items that demand multi-step
inference. This is the diagnostic that distinguishes a knowledge
benchmark from a reasoning benchmark, in the formal frame of
\cref{sec:knowledge-mmlu}: knowledge-dominant $\rho$ has a
zero \glsterm{chain-of-thought}{CoT}-vs-direct gap; reasoning-dominant $\rho$ has a positive
one. (Returning to canonical form: the verifier
$V(\hat{a}, a^\star)$ is the same in both cases; what changes is
the distribution of $\rho$.)
\end{intuition}

\subsection{GPQA Diamond: the Google-proof pivot}
\label{sec:knowledge-gpqa}

\citet{rein2023-gpqa} pivot from broad coverage to narrow-domain depth%
\footnote{KB excerpt: \texttt{kb/excerpts/rein2023-gpqa.md\#abstract}.}
\citep[\S{}abstract]{rein2023-gpqa}:

\begin{quote}\itshape
We present GPQA, a challenging dataset of 448 multiple-choice questions
written by domain experts in biology, physics, and chemistry.
\end{quote}

Construction emphasizes \emph{Google-proofness}: the methodological
contribution is operationalizing what it means for a question to
require expertise that simple web search cannot supply
\citep[\S{}abstract]{rein2023-gpqa}:

\begin{quote}\itshape
highly skilled non-expert validators only reach 34\% accuracy, despite
spending on average over 30 minutes with unrestricted access to the web.
\end{quote}

The $34\%$ non-expert-with-web baseline is the load-bearing
methodological number. It is meaningfully above the $25\%$ random
baseline (web access does help), but well below the $\sim 65\%$
expert-PhD baseline. The gap between $34\%$ and $65\%$ is the operational
definition of ``expertise required even with web search'' --- an
attempt to construct items that resist the contamination failure mode
\cref{sec:knowledge-mmlu} cannot defend against, by making the answer
unrecoverable from public web text alone.

\glsterm{gpqa-diamond-2}{GPQA Diamond}, the canonical leaderboard slice, is the $198$-question
subset where two expert annotators agreed on the answer and at least
one non-expert validator with web access got it wrong
\citep[\S{}sec-diamond]{rein2023-gpqa}. At launch, the strongest
GPT-4-based baseline reached $39\%$, just above the non-expert ceiling
\citep[\S{}sec-ai-baseline]{rein2023-gpqa}. By 2026, frontier reasoning
models are reported above $90\%$ on Diamond --- $\sim 25$ points above
the human PhD baseline. The benchmark is now effectively saturated at
the frontier, in roughly $2.5$ years from publication. The intended
downstream use, per the abstract, was \glsterm{scalable-oversight}{scalable oversight} research:
items so hard that even skilled human supervisors cannot reliably
verify the model's answer. That use case has not yet expired ---
the benchmark continues to function as a \glsterm{probe-2}{probe} for whether weak
supervisors can verify strong models, even where simple
capability-tracking has saturated.

\subsection{FrontierMath: the unpublished-problem benchmark}
\label{sec:knowledge-frontiermath}

\citet{glazer2024-frontiermath} take the contamination defense further.
The benchmark is constructed from \emph{unpublished} problems with
single-canonical-answer formats designed for automated verification%
\footnote{KB excerpt: \texttt{kb/excerpts/glazer2024-frontiermath.md\#abstract}.}
\citep[\S{}abstract]{glazer2024-frontiermath}:

\begin{quote}\itshape
We introduce \glsterm{frontiermath}{FrontierMath}, a benchmark of hundreds of original,
exceptionally challenging mathematics problems crafted and vetted by
expert mathematicians. The questions cover most major branches of
modern mathematics --- from computationally intensive problems in
number theory and real analysis to abstract questions in algebraic
geometry and category theory.
\end{quote}

The difficulty calibration is explicit: ``Solving a typical problem
requires multiple hours of effort from a researcher in the relevant
branch of mathematics, and for the upper end questions, multiple days''
\citep[\S{}sec-difficulty]{glazer2024-frontiermath}. Problems are
authored by working mathematicians (verified PhD or equivalent) and
cross-vetted by a second mathematician for solvability and uniqueness
of answer; problems that do not admit a single canonical answer
(integer, polynomial, or parseable closed-form expression) are
excluded \citep[\S{}sec-verification]{glazer2024-frontiermath}. The
single-canonical-answer constraint is what enables automated grading
without \glsterm{llm-as-judge}{LLM-as-judge} --- and without exposing intermediate work to
gold-data leakage.

The headline launch number is the load-bearing one. As of late 2024,
state-of-the-art models that scored $80\text{--}95\%$ on \glsterm{mmlu}{MMLU} scored
\emph{under $2\%$} on \glsterm{frontiermath}{FrontierMath} \citep[\S{}sec-launch-result,\S{}abstract]{glazer2024-frontiermath}:

\begin{quote}\itshape
Current state-of-the-art AI models solve under 2\% of problems,
revealing a vast gap between AI capabilities and the prowess of the
mathematical community.
\end{quote}

The two-percent floor restored a usable signal for measuring
research-grade math reasoning at exactly the moment AIME and MATH-500
were saturating under reasoning-RL. The benchmark's design choice ---
unpublished items, single-canonical-answer verifiers, multi-day
expert calibration --- is the most aggressive contamination defense in
the lineage and the closest analogue to a closed-book exam where the
question pool itself is private.

\subsection{Humanity's Last Exam: the post-MMLU broad-coverage frontier}
\label{sec:knowledge-hle}

\citet{phan2025-hle} (\glsterm{hle}{Humanity's Last Exam}, \glsterm{hle}{HLE}) is the most recent
attempt to rebuild a broad-coverage knowledge benchmark with frontier
headroom%
\footnote{KB excerpt: \texttt{kb/excerpts/phan2025-hle.md\#abstract}.}.
The motivation paragraph names the gap directly
\citep[\S{}abstract]{phan2025-hle}:

\begin{quote}\itshape
LLMs now achieve over 90\% accuracy on popular benchmarks like \glsterm{mmlu}{MMLU},
limiting informed measurement of state-of-the-art \glsterm{llm}{LLM} capabilities.
\end{quote}

The construction is $2{,}500$ questions across dozens of subjects ---
mathematics, humanities, and the natural sciences ---
positioned as ``the final closed-ended academic benchmark of its kind
with broad subject coverage'' \citep[\S{}abstract]{phan2025-hle}. The
verifier discipline is a deliberate inheritance from
\cref{sec:knowledge-gpqa,sec:knowledge-frontiermath}
\citep[\S{}sec-format]{phan2025-hle}:

\begin{quote}\itshape
a known solution that is unambiguous and easily verifiable, but cannot
be quickly answered via internet retrieval.
\end{quote}

What makes \glsterm{hle}{HLE} distinct from \glsterm{mmlu}{MMLU} and GPQA: it combines multiple-choice
and short-answer formats, includes multimodal items (images) for some
questions, and uses choice-equality or string-equality-after-normalization
verifiers --- \emph{not} \glsterm{llm-as-judge}{LLM-as-judge} --- so the gold standard is
mechanically reproducible. What makes it distinct from \glsterm{frontiermath}{FrontierMath}:
broad subject coverage rather than a single discipline, and item
publication so the test set is auditable. The capability gap at
publication is large --- ``State-of-the-art LLMs demonstrate low
accuracy and calibration on \glsterm{hle}{HLE}, highlighting a significant gap between
current \glsterm{llm}{LLM} capabilities and the expert human frontier on closed-ended
academic questions'' \citep[\S{}sec-capability-gap]{phan2025-hle} ---
which is the property the design was reaching for: a knowledge
benchmark whose ceiling is far above 2026 frontier capability, on the
implicit bet that a benchmark with measurable headroom for $2$--$5$ years
is more valuable than one that saturates in $12$--$24$ months.

\subsection{What the lineage tells us}
\label{sec:knowledge-summary}

The \glsterm{mmlu}{MMLU}-to-\glsterm{hle}{HLE} arc is a coherent progression with a single underlying
move: each successor responds to a specific failure mode of its
predecessor. \glsterm{mmlu}{MMLU} saturated, so \glsterm{mmlu-redux-2}{MMLU-Redux} audited the gold labels
and exposed the $6.49\%$ verifier-error ceiling. The error ceiling
made within-band scores meaningless, so \glsterm{mmlu-pro}{MMLU-Pro} raised difficulty,
widened the choice set, and added the $24$-prompt-style robustness
protocol. Pure broad coverage was no longer the bottleneck, so GPQA
pivoted to narrow-domain depth with a \glsterm{google-proof-construction}{Google-proof construction}. The
contamination defense in GPQA was item-level, so \glsterm{frontiermath}{FrontierMath} went
further with unpublished items and a single-canonical-answer verifier.
Single-discipline depth was insufficient for cross-domain measurement,
so \glsterm{hle}{HLE} rebuilt the broad-coverage frame with the contamination and
verifier discipline of its predecessors carried forward. Each step
buys $12$--$24$ months of frontier headroom; none has produced a
benchmark whose half-life looks fundamentally different from \glsterm{mmlu}{MMLU}'s.

The honest 2026 picture is that knowledge benchmarks at the frontier
are no longer about \emph{what facts the model knows} --- on
undergraduate-level material, scale and pre-training have closed
that gap inside the verifier-error band --- but about whether the
question pool can be kept ahead of contamination, whether the verifier
is right, and whether the prompt format leaks more signal than the
question content. The next layer of the layered defense, taken up in
\cref{sec:contamination}, is the one that asks the prior question
explicitly: \emph{what does it mean for an item to have leaked into
the training distribution}, and how do we tell?
